% Rapatrié depuis le doc de travail (onglet Article FR), révision 144.

\section{Conclusion}
\label{sec:8}

Cet article a proposé une démarche en trois blocs pour étudier et piloter la résilience d'une chaîne logistique à partir de crises simulées. Un générateur soumet un réseau multi-échelon à des perturbations maîtrisées et rejoue chaque crise sous plusieurs politiques de gestion ; une courbe hyperbolique composée caractérise chaque trajectoire par des paramètres de mécanisme et des indicateurs de résilience ; un réseau bayésien causal, dont la structure respecte l'ordre temporel d'une crise et traite la politique comme une intervention, en tire pronostic, diagnostic et préconisation.

Appliquée au réseau ISOMORPH sur 2 000 crises, la démarche montre que la réaction des gestionnaires agit sur la durée d'une crise plutôt que sur sa profondeur : la politique la plus réactive évite 73 \% du service perdu, contre 18 \% et 14 \% pour les politiques prudente et tardive, une hiérarchie retrouvée sur un second réseau. L'efficacité de chaque politique dépend de la nature du choc et de la topologie, et la couverture de stock protège contre les ruptures de distribution mais pas contre les chocs de demande. Le réseau bayésien appris pronostique la classe de perte de 65 \% des crises nouvelles et préconise, lorsque la durée du choc est connue, une politique plus sobre dans environ une crise sur sept pour une perte de réussite inférieure à deux points.

La démarche est ouverte et transposable : tout réseau multi-échelon décrit dans le format proposé peut être étudié de la même façon. Trois prolongements en découlent directement : valoriser le coût des leviers pour arbitrer explicitement entre coût et résilience, imposer les leviers eux-mêmes pour en identifier l'effet causal propre, et confronter les préconisations à des crises réelles.
